\documentclass[11pt,a4paper]{article}
\usepackage[utf8]{inputenc}
\usepackage[T1]{fontenc}
\usepackage[margin=1in]{geometry}
\usepackage{xcolor}
\usepackage{hyperref}
\definecolor{primarycolor}{RGB}{33,37,41}
\definecolor{secondarycolor}{RGB}{0,102,204}
\hypersetup{colorlinks=true,urlcolor=secondarycolor}
\pagestyle{empty}
\setlength{\parindent}{0pt}
\setlength{\parskip}{0.8em}
\begin{document}
Dear Hiring Team,

I am applying for the Cloud Infrastructure / DevOps Contractor role. At Odrava Health Software, I own GitLab CI pipelines for three product teams and built merge-request test environments deployed with Helm to a shared Kubernetes cluster. The opportunity to help develop infrastructure for Corvale’s new train fleet is an appealing way to apply this platform work in a new domain.

I reduced average pipeline time from 38 to 14 minutes and monthly CI runner cost by approximately 30\% through parallelisation, dependency caching and pipeline restructuring. I also automated short-lived environments with PostgreSQL and seeded data, and maintain an approximately 900-test end-to-end suite whose flaky rate fell from 7\% to under 1\%. I documented how our CI works for onboarding.

My Kubernetes work is at the workload level: Odrava’s platform team operates the shared cluster, while I have namespace-level access and build Helm charts for test environments. My personal homelab uses k3s, Argo CD, Prometheus, Grafana and Loki; I have also provisioned a small Hetzner VPS with Terraform. I would welcome the chance to discuss how this experience in CI/CD, automation and observability could contribute to Corvale’s fleet connectivity work.

Sincerely,\\
Marta Kowal-Brandt
\end{document}